\section{Desarrollo del taller}

\subsection{Diagnóstico inicial del caso}

\scenario{1}{Análisis dirigido de MediSalud · fuente literal}

La revisión del caso identifica tres procesos de máxima criticidad. Primero, la atención médica y el registro HCE, porque cualquier demora o pérdida de información afecta la continuidad clínica. Segundo, la prescripción y dispensación, donde un error puede causar daño directo. Tercero, la admisión y facturación, que constituyen la entrada del paciente y sostienen la operación financiera.

Los perfiles más expuestos son el médico, que necesita información completa durante una consulta; el paciente, que agenda, consulta resultados y usa telemedicina; y el personal de admisión, que resuelve disponibilidad, cobertura y cobro. Enfermería y farmacia también requieren sincronización oportuna con la HCE. Gerencia y calidad necesitan una vista agregada, pero no sustituyen la experiencia de los perfiles operativos.

\begin{table}[H]
\centering
\caption{Matriz de análisis inicial}
\begin{tabularx}{\textwidth}{p{3.4cm}Y}
\toprule
Pregunta guía & Respuesta consolidada \\
\midrule
Procesos críticos & Atención y HCE; prescripción y dispensación; agenda, admisión y facturación. \\
Usuarios afectados & Médicos, pacientes y admisión como grupos principales; enfermería y farmacia por dependencia clínica. \\
Evidencia disponible & Incidentes con fecha, módulo, descripción, rol y sede; contexto tecnológico y requisitos no funcionales. \\
Evidencia faltante & Resultados de tarea, duraciones, encuestas, escenarios de red/dispositivo y criterios uniformes de decisión. \\
Decisión metodológica & Preservar incidentes y complementar la observación con datos simulados identificados. \\
\bottomrule
\end{tabularx}
\end{table}

\subsubsection{Preguntas de discusión}

\textbf{1. ¿Por qué la disponibilidad de servidores no demuestra por sí sola una buena calidad en uso?}

Porque un sistema puede permanecer disponible y, aun así, impedir que la persona complete su tarea, exigir tiempos excesivos, producir resultados inseguros o generar una experiencia insatisfactoria. La disponibilidad describe continuidad técnica; la calidad en uso exige observar objetivos, recursos, percepción, riesgos y contexto.

\textbf{2. ¿Qué diferencia existe entre calidad interna, externa y en uso?}

La calidad interna examina estáticamente código y arquitectura; la externa observa el comportamiento del producto ejecutado bajo pruebas controladas; y la calidad en uso mide lo que usuarios concretos consiguen al realizar tareas reales dentro de un contexto definido.

\textbf{3. ¿Qué stakeholder se beneficiaría más del programa de medición?}

El paciente obtiene el beneficio final más importante porque una atención más efectiva, oportuna y segura protege su salud. Sin embargo, el programa debe medir también al médico, cuya interacción con la HCE condiciona directamente la atención, y proporcionar a gerencia evidencia para decidir. Por ello se prioriza al paciente sin aislarlo de los demás perfiles.

\subsubsection{Conclusión parcial}

El diagnóstico confirma que los incidentes disponibles no bastan para explicar la experiencia completa. Se requiere complementar esa evidencia con resultados de tareas, tiempos, percepción y contexto, manteniendo una relación trazable entre cada problema y su efecto sobre usuarios y negocio.

\subsection{Clasificación según ISO/IEC 25022}

\scenario{2}{Clasificación de 3.000 incidentes · fuente literal}

La clasificación de la actividad asignó 1.849 casos a efectividad, 347 a eficiencia, 103 a satisfacción, 361 a libertad de riesgo y 340 a cobertura del contexto. Efectividad concentra el 61,63\% de los registros; sin embargo, la frecuencia no equivale a prioridad. Los 361 casos de libertad de riesgo requieren evaluación inmediata por la gravedad de sus consecuencias potenciales.

\begin{figure}[H]
\centering
\begin{tikzpicture}[x=0.0038cm,y=0.72cm]
  \foreach \label/\value/\y/\barcolor in {Efectividad/1849/5/mediteal,Eficiencia/347/4/mediblue,Satisfacción/103/3/medigold,Libertad de riesgo/361/2/medicoral,Cobertura del contexto/340/1/medigrey}{
    \fill[\barcolor!18] (0,\y-.22) rectangle (2000,\y+.22);
    \fill[\barcolor] (0,\y-.22) rectangle (\value,\y+.22);
    \node[anchor=east,font=\small] at (-25,\y) {\label};
    \node[anchor=west,font=\small\bfseries,text=\barcolor] at (\value+25,\y) {\value};
  }
  \draw[medigrey!50] (0,.45) -- (0,5.55);
\end{tikzpicture}
\caption{Distribución de incidentes por característica de calidad en uso.}
\label{fig:classification}
\end{figure}

\begin{table}[H]
\centering
\caption{Criterio y ejemplo de clasificación}
\begin{tabularx}{\textwidth}{p{3.2cm}YY}
\toprule
Característica & Pregunta de decisión & Ejemplo de efecto \\
\midrule
Efectividad & ¿La tarea no se completa o su resultado es incorrecto? & Una cita no queda confirmada pese a finalizar el formulario. \\
Eficiencia & ¿El resultado se alcanza con tiempo o pasos excesivos? & El guardado de la nota supera ocho segundos. \\
Satisfacción & ¿La experiencia reduce confianza, agrado o disposición de uso? & El paciente abandona una encuesta o reporta confusión. \\
Libertad de riesgo & ¿Existe daño clínico, financiero o de privacidad? & Una dosis incorrecta, factura duplicada o dato de otro paciente. \\
Cobertura del contexto & ¿El fallo depende de sede, dispositivo o conectividad? & Una teleconsulta cae bajo una red doméstica limitada. \\
\bottomrule
\end{tabularx}
\end{table}

La justificación se almacena junto a cada fila procesada. Esta práctica es necesaria porque una descripción puede admitir más de una lectura. Una receta que no se guarda afecta efectividad; una receta guardada con dosis errónea se clasifica principalmente como libertad de riesgo. La característica principal expresa el efecto dominante, aunque el análisis pueda reconocer impactos secundarios.

\begin{longtable}{p{0.9cm}p{2.7cm}p{8.4cm}}
\caption{Clasificación de los incidentes propuestos en la actividad}\label{tab:incidentes-seis}\\
\toprule
ID & Característica principal & Justificación técnica \\
\midrule
\endfirsthead
\toprule
ID & Característica principal & Justificación técnica \\
\midrule
\endhead
1001 & Eficiencia & La nota sí puede guardarse, pero consume 22 segundos; el problema dominante es el tiempo utilizado para alcanzar el resultado. \\
1002 & Efectividad & Después de tres intentos el paciente no consigue agendar, por lo que el objetivo de la tarea no se alcanza. \\
1003 & Libertad de riesgo & La factura duplicada puede ocasionar cobros indebidos y daño financiero, aunque el proceso haya producido una factura. \\
1004 & Cobertura del contexto & La interrupción de la videollamada evidencia que la tarea no se sostiene en el contexto tecnológico de la teleconsulta remota. \\
1005 & Libertad de riesgo & Mostrar datos de otro paciente expone información clínica confidencial y crea un riesgo de privacidad y de decisión médica. Se prioriza el daño potencial sobre el fallo funcional. \\
1006 & Satisfacción & La confusión y el abandono expresan una experiencia negativa y reducen la disposición del paciente para continuar usando el portal. \\
\bottomrule
\end{longtable}

\evidence{La Tabla 2.2 completada se conserva en
\texttt{data/clasificacion/}\allowbreak\texttt{incidentes\_clasificados.csv}.
Este archivo clasifica los 3.000 incidentes del dataset ampliado; las filas 1001--1006 corresponden a los casos mostrados en la plantilla de la actividad.}

\subsubsection{Preguntas de discusión}

\textbf{1. ¿Puede un sistema ser efectivo pero no eficiente?}

Sí. El médico puede guardar correctamente una nota clínica y completar su objetivo, pero tardar 22 segundos. La tarea es efectiva porque obtiene el resultado esperado, aunque no es eficiente porque consume más tiempo que la meta de ocho segundos.

\textbf{2. ¿Por qué la cobertura del contexto es relevante para una red con cinco sedes?}

Porque una función útil en Quito puede degradarse en Manta, Ambato, Cuenca o Guayaquil debido a conectividad, equipos, carga o condiciones operativas diferentes. Evaluar una sola sede no demuestra que el HIS mantenga su calidad para toda la población prevista.

\subsubsection{Conclusión parcial}

La clasificación convierte reportes ambiguos en evidencia comparable. Distinguir el resultado, los recursos, la percepción, el riesgo y el contexto evita concentrar todos los problemas en efectividad y permite asignar prioridades coherentes con su impacto.

\evidence{La salida completa se encuentra en \texttt{apps/data/processed/}\allowbreak\texttt{incidentes\_clasificados.csv}; el resumen de conteos está en \texttt{reportes/evidencias/resultados/}\allowbreak\texttt{clasificacion\_resumen.json}.}

\subsection{Comprensión del modelo SQuaRE}

\scenario{3}{Diferenciación de los niveles de calidad · fuente literal}

\subsubsection{Investigación dirigida}

La \textbf{calidad interna} observa el código y la arquitectura sin ejecutar el producto: complejidad, duplicación, mantenibilidad y estructura. La \textbf{calidad externa} estudia el comportamiento del sistema ejecutado en un entorno controlado: respuesta, estabilidad y errores bajo una carga definida. Ambas perspectivas se relacionan con el modelo de calidad de ISO/IEC 25010 \cite{iso25010}.

La \textbf{calidad en uso} observa lo que consigue una persona al realizar una tarea dentro de un contexto. Considera efectividad, eficiencia, satisfacción, libertad de riesgo y cobertura del contexto; ISO/IEC 25022 proporciona medidas para esta perspectiva \cite{iso25022}. ISO/IEC 25040 conecta las evidencias y medidas con un proceso de evaluación y decisión \cite{iso25040}.

En esta relación, ISO/IEC 25010 define \emph{qué} características conforman el modelo de calidad, mientras ISO/IEC 25022 establece \emph{cómo} medir la calidad en uso mediante medidas, fórmulas y escalas aplicadas a usuarios, tareas y contextos.

\subsubsection{Aplicación al caso MediSalud}

\begin{figure}[H]
\centering
\resizebox{\textwidth}{!}{\begin{tikzpicture}[
  font=\sffamily,
  >={Latex[length=2.6mm]},
  norm/.style={rounded corners=2.5mm,minimum width=4.55cm,minimum height=1.45cm,align=center,font=\sffamily\bfseries,line width=1pt},
  branch/.style={rounded corners=2.5mm,text width=4.2cm,minimum height=5.75cm,align=left,inner sep=8pt,line width=.9pt},
  connector/.style={-{Latex[length=2.6mm]},line width=1.1pt}
]
  \node[rounded corners=3mm,fill=medidark,text=white,minimum width=7.5cm,minimum height=1.55cm,align=center] (square) at (0,13.35) {\Large\bfseries ISO/IEC 25000\\[-1mm]
    \small\normalfont Familia SQuaRE: requisitos, modelos, medición y evaluación};

  \node[norm,draw=mediblue,fill=mediblue!8,text=medidark] (n25010) at (-5.2,10.65)
    {ISO/IEC 25010\\[-1mm]\small\normalfont Modelo de calidad};
  \node[norm,draw=mediteal,fill=mediteal!8,text=medidark] (n25022) at (0,10.65)
    {ISO/IEC 25022\\[-1mm]\small\normalfont Medición de calidad en uso};
  \node[norm,draw=medicoral,fill=medicoral!8,text=medidark] (n25040) at (5.2,10.65)
    {ISO/IEC 25040\\[-1mm]\small\normalfont Proceso de evaluación};

  \draw[connector,draw=mediblue] (square.south) -- ++(0,-.48) -| (n25010.north);
  \draw[connector,draw=mediteal] (square.south) -- (n25022.north);
  \draw[connector,draw=medicoral] (square.south) -- ++(0,-.48) -| (n25040.north);

  \node[branch,draw=mediblue!70,fill=mediblue!5] (product) at (-5.2,6.65) {
    \textcolor{mediblue}{\fontsize{9.5}{11}\selectfont\bfseries PRODUCTO SOFTWARE}\\[-1mm]
    {\scriptsize ISO/IEC 25010}\par\vspace{2.5mm}
    \begin{minipage}{3.95cm}
      \footnotesize
      \textbf{Calidad interna}\\
      \textcolor{mediblue}{Vista estática}\\
      Código y arquitectura sin ejecutar.\\
      \textit{SonarQube: complejidad 8.}\\[2.5mm]
      \color{mediblue!45}\rule{\linewidth}{.5pt}\color{black}\\[-1mm]
      \textbf{Calidad externa}\\
      \textcolor{medigold!80!black}{Vista dinámica controlada}\\
      Producto ejecutado bajo carga.\\
      \textit{JMeter: 500 usuarios, 0\% de error.}
    \end{minipage}
  };

  \node[branch,draw=mediteal!70,fill=mediteal!5] (use) at (0,6.65) {
    \textcolor{mediteal}{\fontsize{9.5}{11}\selectfont\bfseries EXPERIENCIA DE USO}\\[-1mm]
    {\scriptsize ISO/IEC 25022}\par\vspace{2.5mm}
    \begin{minipage}{3.95cm}
      \footnotesize
      \textbf{Calidad en uso}\\
      Usuario + objetivo + tarea + recursos + entorno.\\[2.5mm]
      \colorbox{mediteal!14}{\textcolor{mediteal}{\scriptsize\bfseries Efectividad}}\quad
      \colorbox{mediblue!13}{\textcolor{mediblue}{\scriptsize\bfseries Eficiencia}}\\[1.5mm]
      \colorbox{medigold!18}{\textcolor{medigold!70!black}{\scriptsize\bfseries Satisfacción}}\\[1.5mm]
      \colorbox{medicoral!13}{\textcolor{medicoral}{\scriptsize\bfseries Libertad de riesgo}}\\[1.5mm]
      \colorbox{medigrey!13}{\textcolor{medigrey}{\scriptsize\bfseries Cobertura del contexto}}\\[2mm]
      \textit{HCE simulada: P90 de 11,71 s; meta de 8 s.}
    \end{minipage}
  };

  \node[branch,draw=medicoral!70,fill=medicoral!5] (evaluation) at (5.2,6.65) {
    \textcolor{medicoral}{\fontsize{9.5}{11}\selectfont\bfseries DECISIÓN SISTEMÁTICA}\\[-1mm]
    {\scriptsize ISO/IEC 25040}\par\vspace{2.5mm}
    \begin{minipage}{3.95cm}
      \footnotesize
      \textbf{1.} Establecer requisitos\\[1.5mm]
      \textbf{2.} Especificar la evaluación\\[1.5mm]
      \textbf{3.} Diseñar la medición\\[1.5mm]
      \textbf{4.} Ejecutar y comparar\\[1.5mm]
      \textbf{5.} Concluir y recomendar\\[2.5mm]
      Evidencia + criterio = decisión trazable.
    \end{minipage}
  };

  \draw[connector,draw=mediblue] (n25010.south) -- (product.north);
  \draw[connector,draw=mediteal] (n25022.south) -- (use.north);
  \draw[connector,draw=medicoral] (n25040.south) -- (evaluation.north);

  \node[rounded corners=2mm,fill=medidark,text=white,text width=14.65cm,minimum height=1.1cm,align=center,font=\small] at (0,2.35)
    {La calidad interna facilita un producto mantenible; la calidad externa comprueba su comportamiento; la calidad en uso determina si las personas alcanzan sus objetivos en condiciones concretas.};
\end{tikzpicture}
}
\caption{Mapa conceptual de la familia SQuaRE aplicado a MediSalud HIS.}
\label{fig:square}
\end{figure}

\evidence{El mapa también se entrega como archivo independiente en
\texttt{reportes/mapa-conceptual/}\allowbreak\texttt{mapa\_conceptual\_square.pdf}
y como vista previa PNG en la misma carpeta y en \texttt{docs/mapa-conceptual}.}

\begin{table}[H]
\centering
\caption{Los tres niveles de calidad aplicados a MediSalud HIS}
\begin{tabularx}{\textwidth}{p{2.4cm}YY}
\toprule
Nivel & Actividad ejecutada & Evidencia obtenida \\
\midrule
Calidad interna & SonarQube Community Build sobre el servicio de facturación. & Complejidad ciclomática 8; cognitiva 3; 54 líneas; 0 bugs y 0 vulnerabilidades. \\
Calidad externa & JMeter con 500 usuarios sincronizados consultando citas por el Gateway. & 500 respuestas correctas; 0\% de error; promedio 1.602 ms; P90 2.002 ms. \\
Calidad en uso & Medición del registro HCE durante la simulación de consulta externa. & P90 de 11,71 s frente a la meta de 8 s; resultado amarillo. \\
\bottomrule
\end{tabularx}
\end{table}

\evidence{Las respuestas reproducibles se conservan en
\texttt{reportes/evidencias/calidad/}\allowbreak\texttt{sonarqube/}\allowbreak\texttt{metricas.json} y
\texttt{reportes/evidencias/calidad/}\allowbreak\texttt{jmeter/}\allowbreak\texttt{resumen.json}.
Los resultados JTL son regenerables y permanecen fuera del control de versiones.}

\subsubsection{Preguntas de discusión}

\textbf{1. ¿Puede un sistema tener excelente calidad interna y mala calidad en uso?}

Sí. Un módulo puede presentar baja complejidad, ausencia de problemas detectados y una arquitectura ordenada, pero aun así obligar al usuario a esperar o repetir acciones. En MediSalud, SonarQube no detectó bugs ni vulnerabilidades en facturación y reportó complejidad ciclomática 8; sin embargo, la calidad en uso del conjunto conserva un P90 de 11,71 segundos al registrar HCE y un éxito de agendamiento de 87,86\%. La estructura del código no garantiza que la tarea resulte rápida, completa o satisfactoria.

\textbf{2. ¿Por qué SonarQube no es suficiente para resolver la lentitud percibida por los médicos?}

Porque SonarQube analiza propiedades estáticas y defectos del código, pero no reproduce la consulta clínica ni mide tiempo de red, base de datos, concurrencia, integración o esfuerzo del usuario. Para comprender la lentitud se necesitan pruebas dinámicas, trazas operativas y medidas de calidad en uso. JMeter aporta comportamiento externo bajo carga e ISO/IEC 25022 aporta la relación entre usuario, tarea, contexto y resultado.

\subsubsection{Conclusión parcial}

La calidad en uso es la perspectiva más próxima al negocio y al paciente, pero se sostiene sobre la calidad interna y externa. SonarQube ayuda a conocer el estado estático del producto; JMeter muestra su comportamiento controlado; y las medidas ISO/IEC 25022 determinan si las personas logran sus objetivos en condiciones de uso. Ninguna de estas vistas sustituye a las demás.

\subsection{Atributos, tareas y prioridades}

\scenario{4--5}{Atributos UTC y mapeo tarea--característica · reconstrucción declarada}

Se seleccionaron tareas observables que cubren los perfiles del caso y pueden producir evidencia. La prioridad combina impacto clínico, financiero, operativo y reputacional con la urgencia sugerida por los incidentes.

\begin{longtable}{p{2.0cm}p{3.2cm}p{3.7cm}p{2.7cm}p{1.6cm}}
\caption{Catálogo de tareas, contextos y atributos}\label{tab:utc}\\
\toprule
Usuario & Tarea & Contexto & Atributo observable & Prioridad \\
\midrule
\endfirsthead
\toprule
Usuario & Tarea & Contexto & Atributo observable & Prioridad \\
\midrule
\endhead
Médico & Guardar nota HCE & Consulta, hora pico & P90 de duración & Crítica \\
Médico & Revisar alergias & Antes de prescribir & Disponibilidad correcta & Crítica \\
Paciente & Agendar cita & Móvil y 4G & Tasa de éxito & Alta \\
Paciente & Realizar teleconsulta & WiFi doméstico & Tasa de caídas & Alta \\
Admisión & Emitir factura & Seguro y cierre mensual & Tasa de error & Crítica \\
Farmacia & Validar receta & Turno nocturno & Recetas correctas & Crítica \\
Enfermería & Registrar signos & Tableta y red hospitalaria & Sincronización & Alta \\
Médico & Cargar imagen & Dispositivo clínico & P90 de carga & Alta \\
Paciente & Responder encuesta & Aplicación móvil & Abandono y CSAT & Media \\
Gerencia & Revisar KPI & Consolidado multisede & Cobertura de datos & Alta \\
\bottomrule
\end{longtable}

Las tareas críticas no son necesariamente las más frecuentes. La prescripción se prioriza por severidad; la HCE por su relación con la decisión clínica; y la facturación por el riesgo de cobros incorrectos. El portal de citas y la telemedicina tienen prioridad alta porque son puntos de acceso directo del paciente y dependen de contextos heterogéneos.

\subsection{Diseño de métricas}

\scenario{6}{Fichas de medida y criterios de decisión · reconstrucción declarada}

Cada ficha incluye código, nombre, característica, objetivo, fórmula, unidad, fuente, periodicidad, meta y semáforo. El valor no se interpreta sin su meta: 96,52\% puede parecer alto para una función general, pero es rojo cuando el resultado restante representa recetas incorrectas.

Las fórmulas principales son:

\begin{align}
\text{Éxito}(T) &= \frac{\text{tareas exitosas}}{\text{tareas iniciadas}}\times 100, \\
\text{Error}(T) &= \frac{\text{tareas no exitosas}}{\text{tareas iniciadas}}\times 100, \\
\text{Abandono} &= \frac{\text{encuestas no completadas}}{\text{encuestas iniciadas}}\times 100, \\
\text{CSAT} &= \frac{\sum_{i=1}^{n} \text{puntuación}_{i}}{n}, \\
\text{Variabilidad} &= \sigma\left(\overline{t}_{\text{sede }1},\ldots,\overline{t}_{\text{sede }5}\right).
\end{align}

El percentil 90 se define como el menor valor que contiene al 90\% de las observaciones ordenadas. Se eligió para HCE e imagenología porque hace visible la experiencia lenta del extremo superior sin depender de un único máximo atípico.

\begin{longtable}{p{1.6cm}p{3.6cm}p{2.6cm}p{2.0cm}p{2.1cm}}
\caption{Catálogo resumido de métricas}\label{tab:catalog}\\
\toprule
Código & Métrica & Característica & Unidad & Meta \\
\midrule
\endfirsthead
\toprule
Código & Métrica & Característica & Unidad & Meta \\
\midrule
\endhead
M-EFI-01 & P90 registro HCE & Eficiencia & segundos & $\leq 8$ \\
M-EFE-01 & Éxito de agendamiento & Efectividad & porcentaje & $\geq 95$ \\
M-RIE-01 & Errores de facturación & Libertad de riesgo & porcentaje & $<1$ \\
M-SAT-01 & Abandono de encuesta móvil & Satisfacción & porcentaje & $<10$ \\
M-CC-01 & Variabilidad HCE entre sedes & Cobertura del contexto & segundos & $\leq 2$ \\
M-RIE-02 & Incidentes de privacidad & Libertad de riesgo & incidentes & $=0$ \\
M-SAT-02 & CSAT promedio & Satisfacción & escala 1--5 & $\geq 4$ \\
M-EFE-02 & Recetas correctas & Efectividad & porcentaje & $=100$ \\
M-EFI-02 & P90 carga de imagen & Eficiencia & segundos & $\leq 10$ \\
M-CC-02 & Caídas de teleconsulta & Cobertura del contexto & porcentaje & $<5$ \\
\bottomrule
\end{longtable}

\subsection{Datos, automatización y KPI}

\scenario{7--9}{Obtención, procesamiento y visualización · reconstrucción declarada}

La simulación instrumenta cada flujo con un identificador de auditoría y un escenario. En operación normal se registra la finalización esperada; en modo defectuoso se activa una regla controlada como demora, pérdida de guardado, pasos adicionales, doble facturación, dosis incompatible o desconexión. La probabilidad y la semilla permiten repetir la distribución sin convertir el fallo en una vulnerabilidad real.

El pipeline analítico ejecuta cuatro responsabilidades separadas:

\begin{enumerate}
  \item \textbf{Clasificación:} carga los incidentes, normaliza campos y asigna característica y justificación.
  \item \textbf{Simulación:} genera eventos y encuestas con periodo, sede, rol, módulo, origen y semilla.
  \item \textbf{Medición:} agrupa observaciones, calcula porcentajes, percentiles, medias y desviación entre sedes.
  \item \textbf{Exportación:} guarda CSV y JSON para que API, tablero e informe consuman la misma evidencia.
\end{enumerate}

La API de medición expone salud, incidentes, métricas, resumen y ejecuciones de simulación. React y Flutter acceden mediante el Gateway. El tablero permite leer valores agregados sin incorporar métricas académicas en las tareas del paciente. Esta separación conserva el carácter realista del HIS: la persona clínica atiende; el paciente gestiona su salud; gerencia observa tendencias.

\begin{table}[H]
\centering
\caption{Relación entre fallo, evidencia y medida}
\begin{tabularx}{\textwidth}{p{3.4cm}YY}
\toprule
Fallo controlado & Evidencia generada & KPI afectado \\
\midrule
Guardado HCE lento & Duración y resultado de \texttt{hce\_save} & M-EFI-01 \\
Agendamiento prolongado & Inicio, pasos y finalización de \texttt{appointment} & M-EFE-01 \\
Factura duplicada & Estado y bandera de duplicidad de \texttt{billing} & M-RIE-01 \\
Dosis incorrecta & Resultado de validación de \texttt{prescription} & M-EFE-02 \\
Imagen lenta & Duración de \texttt{imaging} & M-EFI-02 \\
Teleconsulta interrumpida & Estado de \texttt{teleconsultation} & M-CC-02 \\
\bottomrule
\end{tabularx}
\end{table}

\evidence{La automatización reside en \texttt{apps/analytics}; los datos regenerables en \texttt{apps/data/simulated} y \texttt{apps/data/processed}; y los resúmenes seleccionados en \texttt{reportes/evidencias/resultados}.}
